\documentclass[11pt]{article}
\usepackage[margin=25mm]{geometry}
\usepackage{fontspec}
\setmainfont{Latin Modern Roman}
\usepackage{amsmath,amssymb,mathtools}
\usepackage{unicode-math}
\setmathfont{Latin Modern Math}
\usepackage{xurl,hyperref}
\hypersetup{colorlinks=true,urlcolor=blue}
\setcounter{secnumdepth}{0}
\setlength{\parindent}{0pt}
\setlength{\parskip}{5pt}
\setlength{\emergencystretch}{3em}
\newcommand{\sref}[2]{\hyperlink{source:#1:#2}{[S#2]}}
\newcommand{\cataloguescope}[1]{\paragraph{Recorded target.}#1}
\newcommand{\R}{\mathbb R}
\newcommand{\C}{\mathbb C}
\newcommand{\N}{\mathbb N}
\newcommand{\Z}{\mathbb Z}
\newcommand{\Q}{\mathbb Q}
\newcommand{\F}{\mathbb F}
\newcommand{\eps}{\varepsilon}
\begin{document}
% UnsolvedMath ID 30000244; source catalogue code OWR-832-003
\setcounter{section}{30000243}
\section{Physical Interpretation of Elliptic Partition Functions}\label{30000244}
{\small UnsolvedMath ID 30000244 · Mathematical Physics}
\subsection{Definitions and mathematical statement}

\paragraph{Shared notation.}$\mathbb N=\{1,2,\ldots\}$; the standard number fields have their usual meanings. All additional conventions are stated locally.


\paragraph{Definitions and conventions.} A partition function is the amplitude or generating function assigned to a physical theory on its spacetime background. Elliptic cohomology is a generalized cohomology theory associated with elliptic curves; the report uses it to refine the type IIA and IIB flux partition functions. The proposed F-theory interpretation introduces two extra dimensions represented by an elliptic curve. An index is the signed dimension of the kernel and cokernel of a suitable Dirac-type operator; the $E_8$ and Rarita--Schwinger indices refer to gauge and spin-$3/2$ contributions in the cited construction.

\paragraph{Source scope and qualifications.} The twelve-dimensional physical theory and its loop indices are proposed constructions. The report does not supply a fully axiomatized F-theory partition function with all backgrounds fixed; the target is an interpretation and its proposed tests, not an already defined equality for every conceivable background.

\paragraph{Catalogue mathematical target.}
\paragraph{Catalogue wording retained for traceability.}
Do the elliptic-cohomology partition functions for type II string theory coincide with partition functions of twelve-dimensional F-theory compactified on an elliptic curve? Test this proposed interpretation by computing suitable loop versions of the $E_8$ and Rarita–Schwinger indices.

\paragraph{Saved literature context.} The supplied catalogue's assessment dated 2026-08-21 records: The report states that the physical conjecture had not been verified; later index and partition-function work gives evidence and special calculations, not a general verification. The assessment is retained as dated provenance; it is not a new comprehensive literature or open-status certification.
\subsection{Short English statement}
Test whether elliptic-cohomology string partition functions arise by compactifying a twelve-dimensional F-theory, using the proposed loop-index calculations.
\subsection{Research attempt}
\paragraph{Model, date and investigated scope.}
OpenAI GPT-6 Astra (configured model \texttt{gpt-6-astra}), ultra reasoning; generated 11 October 2026. Ordinary product Dirac indices vanish on a flat elliptic factor with pulled-back twists, showing why that shortcut cannot supply the requested loop-index test.

\paragraph{Sources and literature check.}
The complete supplied definition, including source qualifications and the dated literature assessment, was reviewed. The calculation below is an elementary restricted attempt motivated by [S2]; it does not claim novelty or a fresh literature-status determination.

\paragraph{Context and investigated scope.} The complete supplied definition, including source qualifications and the dated literature assessment, was reviewed. The calculation below is an elementary restricted attempt motivated by [S2]; it does not claim novelty or a fresh literature-status determination.

\paragraph{Testing the ordinary product-index shortcut.} The suggested loop-index test cannot be replaced without justification by an ordinary Dirac index. To see a concrete obstruction, take a closed even-dimensional spin manifold $X$ and a flat two-torus $T^2$, with a bundle pulled back from $X$ and product connections. For the graded product Dirac operator,
\[
 D_{X\times T^2}=D_X\otimes I+\Gamma_X\otimes D_{T^2},\qquad
 D_{X\times T^2}^2=D_X^2\otimes I+I\otimes D_{T^2}^2.
\]
Here $\Gamma_X$ is the chirality operator, which anticommutes with $D_X$. Taking the inner product with a zero mode shows that both nonnegative square terms vanish, so the kernel is the tensor product of the two kernels. Graded dimensions multiply.

\paragraph{The elliptic factor kills the ordinary index.} With the periodic spin structure on the flat torus, its positive and negative chiral spin bundles each have one constant zero mode; every nonzero Fourier mode has strictly positive squared eigenvalue. Their index difference is zero. With a nontrivial flat spin structure there is no zero Fourier mode and again the index is zero. Therefore the twisted product index in this setting is zero regardless of the index on $X$. The same conclusion holds for any finite sum or difference of twists pulled back from $X$.

\paragraph{Why this matters and what is missing.} This is a countercheck on a proposed method, not a disproof of the source's F-theory interpretation: loop indices carry additional graded information which an ordinary index forgets. It also does not define a twelve-dimensional path integral, flux quantization, or the proposed $E_8$ and Rarita--Schwinger loop operators. Defining those operators and showing what replaces the product-kernel calculation is the first gap. Both the partition-function comparison and its requested loop-index tests remain unsettled.

\paragraph{Verification and adversarial audit.}
Checked product anticommutation, positivity of squared operators, periodic and nonperiodic torus spin structures, and the restriction to pulled-back twists. Model self-audit only; no independent human proof checking. Prompts A and B were applied to this canonical record and its qualified source scope. Checked product anticommutation, positivity of squared operators, periodic and nonperiodic torus spin structures, and the restriction to pulled-back twists. Retained unresolved status and zero original-target progress; The relevant loop operators and a common definition of the two partition functions have not been constructed.

Coordinator checked the anticommuting product Dirac square, the kernel tensor decomposition, and zero index in both torus spin cases. The product calculation does not establish the loop-space target. This additional check was model-assisted, not independent human verification.

\paragraph{Tools, process and contributions.}
Python 3 was used for file preparation and structural checks; no numerical experiment or theorem prover was used. Contribution prompt A was followed by prompt B. GPT-6 Astra with ultra reasoning drafted the research and performed its stated self-audit; any additional coordinator contribution is identified in the verification and audit. A GPT-6-based Codex coordinator assembled the manuscript and checked repository integration; its exact variant and reasoning configuration were not exposed. No human mathematical contribution, independent human proof check or formal proof certification is supplied. Complete derivations and the consulted source references are included here; private reasoning transcripts are not published.

\paragraph{Final status and first remaining gap.}
\textbf{UNRESOLVED}. The relevant loop operators and a common definition of the two partition functions have not been constructed.

\section{Completion Estimate}
\noindent\textbf{COMPLETION: 0\%}
The calculation is an auxiliary or restricted attempt and establishes no newly verified portion of the full source target. The relevant loop operators and a common definition of the two partition functions have not been constructed.

\subsection{Sources}
\begin{itemize}
\item[\hypertarget{source:30000244:1}{S1}] ULAM.AI and the UnsolvedMath Contributors, UnsolvedMath, record 30000244 (OWR-832-003). Catalogue curation under CC BY 4.0; definitions and English exposition added locally.
\url{https://huggingface.co/datasets/ulamai/UnsolvedMath}.
\item[\hypertarget{source:30000244:2}{S2}] Geometric Topology and Connections with Quantum Field Theory (2005), Oberwolfach Reports, DOI 10.4171/owr/2005/27, pp. 1505--1546; official report source and mathematical context.
\url{https://doi.org/10.4171/owr/2005/27}.
\item[\hypertarget{source:30000244:3}{S3}] Official EMS Press PDF of the same Oberwolfach report; PDF page 11.
\url{https://ems.press/content/serial-article-files/46001}.
\item[\hypertarget{source:30000244:4}{S4}] Kriz--Sati Type II modularity; reference in the saved dated catalogue literature review.
\url{https://arxiv.org/abs/hep-th/0501060}.
\end{itemize}
\label{end:30000244}
\end{document}
